\documentclass[11pt]{article}
\usepackage{fontspec}
\usepackage[french]{babel}
\usepackage{amsmath}
\usepackage{amssymb}
\usepackage{amsthm}
\usepackage{mathrsfs}
\usepackage{tikz-cd}
\usetikzlibrary{babel}
\usepackage{enumitem}

\newcommand{\printedpage}[1]{\marginpar{\footnotesize p.~#1}}
\newcommand{\msp}[1]{\marginpar{\footnotesize [#1]}}
\newcommand{\Ind}{\operatorname{Ind}}
\newcommand{\uHom}{\underline{\mathrm{Hom}}}
\newcommand{\Ob}{\operatorname{Ob}}
\newcommand{\card}{\operatorname{card}}
\newcommand{\Fl}{\mathrm{Fl}}
\newcommand{\PF}{\mathrm{PF}}

\begin{document}

\printedpage{24}
% head: A. Grothendieck : Dérivateurs, Ch. XVIII, Édition des Universités de Montpellier II et Paris VII
% foot: Version 17 décembre 2025

\textit{avec}
\[
\Psi(i) \;=\; \varinjlim_{j\in J}{}_{C}\,\psi(i,j)\;.
\]
\textit{c)} $\Ind_\pi C$ \textit{est stable par petites limites inductives} (\footnote{rajouter aussi, en d), le corollaire 4.4.16, page 101.}).

Il reste à établir ce dernier point. Comme on a déjà les $\varinjlim$ finies par b, il reste à prouver l'existence des petites $\varinjlim$ filtrantes, ou au choix des sommes infinies, ce qui sera plus joli en fait. Mais
\msp{page 97}
on pourra aussi bien traiter le cas général d'une $\varinjlim$ de type $J$ quelconque ($J$ petite). Pour ceci, soit
\[
I \;=\; I(J) \;=\; \begin{aligned}[t]
&\text{ensemble des sous-catégories } J' \text{ de } J \text{ telles que}\\
&\card\Fl(J')\leqslant\pi,
\end{aligned}
\]
ordonnons $I$ par inclusion. Il est évident que $I\neq\emptyset$ (car $\emptyset\in I$), et même que $I$ \textit{est grand devant} $\pi$, car si $(J'_\alpha)_{\alpha\in A}$ avec $\card(A)\leqslant\pi$ est une famille d'éléments de $I$, alors la sous-catégorie de $J$ engendrée par les $J'_\alpha$ est encore dans $I$ (vérification immédiate), et majore les $J'_\alpha$. En fait, la considération de $I(J)$ associée à $J$ va nous permettre de prouver ceci, qui généralise 4.4.13 c), compte tenu de b) :

\textbf{Corollaire 4.4.14} (\footnote{Comparer aussi 4.11.3, page 270.}). \textit{Soit $\mathcal{M}$ une catégorie, $\pi$ un cardinal infini tel que}
\begin{enumerate}
\item[a)] \textit{$\mathcal{M}$ est}
\msp{page 98}
\textit{stable par $\varinjlim$ de cardinal $\leqslant\pi$, et}
\item[b)] \textit{$\mathcal{M}$ est stable par $\varinjlim$ sur les ensembles ordonnés $I$ grands devant $\pi$.}
\end{enumerate}
\textit{Alors $\mathcal{M}$ est stable par petites $\varinjlim$ .}

(\footnote{\textbf{Corollaire 4.4.14.1.} $\mathcal{M}$ étant comme dans 4.4.14, et $f:\mathcal{M}\longrightarrow\mathcal{E}$ étant un foncteur qui commute aux $\varinjlim$ de cardinal $\leqslant\pi$, et aux petites $\varinjlim$ grandes devant $\pi$, alors $f$ commute aux petites $\varinjlim$ .}).

En effet, si $\varphi:J\longrightarrow\mathcal{M}$ est donné, on en construit
\[
\psi\;:\;I\;=\;I(J)\;\longrightarrow\;\mathcal{M}
\]
par
\[
\psi(J')\;=\;\varinjlim_{J'}\varphi(J')\;,
\]
qui existe par hypothèse sur $\mathcal{M}$. Il est alors immédiat que $\varinjlim\limits_{J}\varphi$ existe si et seulement si $\varinjlim\limits_{I}\psi$ existe, or cette existence est assurée du fait que $I$ est grand devant $\pi$, et l'hypothèse sur $\mathcal{M}$.

Pour achever de prouver le théorème 4.4.13, on est donc ramené au cas particulier des limites inductives $\varinjlim\limits_{J}\varphi(j)$, quand $\varphi:J\longrightarrow\Ind'_\pi C$ est donné, avec $J$ ensemble
\printedpage{25}
% head: A. Grothendieck : Dérivateurs, Ch. XVIII, Édition des Universités de Montpellier II et Paris VII
% foot: Version 17 décembre 2025
\msp{page 99}
ordonné grand devant $\pi$. Mais ceci va résulter du résultat plus précis :

\textbf{Corollaire 4.4.15} (\footnote{comparer à 4.7.3.1, qu'il y aurait lieu d'expliciter ici.}). \textit{Soit $C$ une petite catégorie. Alors la sous-catégorie strictement pleine} $\Ind'_\pi C$ \textit{de $C^\wedge$ est stable dans $C^\wedge$ par $\varinjlim$ grandes devant $\pi$, en d'autres termes, ces $\varinjlim$ existent} dans $\Ind'_\pi(C)$, \textit{et le foncteur d'inclusion y commute.} % ?? uncertain: font of "existent" -- italic (as read) or upright like "dans"

\textsc{Démonstration}. Soit $I$ ensemble ordonné, grand devant $\pi$, et $i\longmapsto F_i$ un foncteur $I\longrightarrow\Ind'_\pi(C)$, $F\;=\;\varinjlim\limits_{I}F_i$ dans $C^\wedge$, prouvons que $F\in\Ob\Ind'_\pi(C)$, i.e. (4.4.3) que $C/F$ est grande devant $\pi$. J'utilise le critère b) de 4.4.4, il faut vérifier $\PF_\pi\,1$, $\PF_\pi\,2$. Soient
\[
x_\alpha\xrightarrow{\;u_\alpha\;}F\;=\;\varinjlim F_i
\]
des objets de $C/F$ indexés par $A$ avec $\card A\leqslant\pi$. Chaque $u_\alpha$ provient de $x_\alpha\longrightarrow F_{i_\alpha}$, et en prenant un majorant commun des $i_\alpha$ ($I$ grand devant $\pi$), on trouve un $i\in I$ et des $v_\alpha:x_\alpha\longrightarrow F_i$. Comme $F_i$ est dans $\Ind'_\pi(C)$, i.e. $C/F_i$ grande devant $\pi$, considérant les $v_\alpha$ comme des objets de $C/F_i$, ils sont majorés par un objet $x\longrightarrow F_i$ de $C/F_i$. On a donc
\[
\begin{tikzcd}
x_\alpha \arrow[r] \arrow[d] & F_{i_\alpha} \arrow[d] \arrow[dr, dotted] & \\
x \arrow[r] & F_i \arrow[r] & F\;,
\end{tikzcd}
\]
qui montre que les $x_\alpha\longrightarrow F$ sont majorés par $x\longrightarrow F$,
\msp{page 100}
ce qui prouve $\PF_\pi\,1$. Prouvons $\PF_\pi\,2$, i.e. soient $x\xrightarrow{\;u\;}F$, $y\xrightarrow{\;v\;}F$ deux objets de $C/F$, et $(f_\alpha)_{\alpha\in A}$ une famille de cardinal $\leqslant\pi$ de flèches de $u$ dans $v$. On sait que $u$, $v$ proviennent de $u_i:x\longrightarrow F_i$, $v_i:y\longrightarrow F_i$ pour $i$ assez grand.
\[
\begin{tikzcd}[column sep=small]
x \arrow[rr, "f_\alpha"] \arrow[dr, "u"'] & & y \arrow[dl, "v"] \\
& F &
\end{tikzcd}
\]
\[
\begin{tikzcd}[column sep=small]
x \arrow[rr, "f_\alpha"] \arrow[dr, "u_i"'] & & y \arrow[dl, "v_i"] \\
& F_i &
\end{tikzcd}
\]

Pour tout $\alpha$, on peut choisir $i_\alpha\geqslant i$ tel que $v_if_\alpha$ et $u_i$ deviennent égaux en composant avec $F_i\longrightarrow F_{i_\alpha}$, et comme $\card A\leqslant\pi$ et $I$ grand devant $\pi$, il existe un majorant $i'$ des $i_\alpha$. Au total, on peut supposer (remplaçant $i$ par $i'$) que $v_if_\alpha=u_i$ pour tout $\alpha$. Alors $u_i$, $v_i$ étant considérés comme objets de $C/F_i$, les $f_\alpha$ sont des flèches de $C/F_i$ de $u_i$ dans $v_i$. Comme $C/F_i$ est grand devant $\pi$ (par l'hypothèse $F_i\in\Ob\Ind'_\pi(C)$), on trouve qu'il existe un égalisateur $(y,v_i)\xrightarrow{\;g\;}(z,w_i)$.
\printedpage{26}
% head: A. Grothendieck : Dérivateurs, Ch. XVIII, Édition des Universités de Montpellier II et Paris VII
% foot: Version 17 décembre 2025
\[
\begin{tikzcd}
x \arrow[r, shift left=2, "(f_\alpha)"] \arrow[r] \arrow[r, shift right=2] \arrow[dr, "u_i"'] & y \arrow[r, "g"] \arrow[d, "v_i"'] & z \arrow[dl, "w_i"] \\
& F_i &
\end{tikzcd}
\]

Alors l'image de $g$ dans $C/F$ (via $F_i\longrightarrow F$) égalise les flèches $f_\alpha$, en tant que flèches de $C/F$, ce qui prouve $\PF_\pi\,2$ et achève la démonstration de 4.4.14, et par là aussi celle du théorème 4.4.13, et du lemme 4.3. Par là est aussi achevée la démonstration de la première assertion du lemme
\msp{page 101}
4.3.7 (p. 77), ce qui achève la démonstration du théorème 4.3.6 (p. 70-71). Il reste à établir la propriété universelle énoncée dans 4.3.7, que je vais rappeler :

\textbf{Corollaire 4.4.16} (\footnote{Pourrait être incorporé en partie d) au théorème 4.4.13, car c'est un complément important, et pas tellement évident finalement.}). \textit{Soient $C$ une catégorie équivalente à une petite catégorie, stable par $\varinjlim$ de cardinal $\leqslant\pi$, où $\pi$ est un cardinal infini donné, et considérons le foncteur d'inclusion canonique}
\[
C\xrightarrow{\;\alpha\;}\Ind_\pi(C)\;. \tag*{(4.4.16.1)}
\]
\textit{Ce foncteur est} 2\textit{-universel pour les foncteurs $f:C\longrightarrow\mathcal{E}$, où $\mathcal{E}$ est une $\mathcal{U}$-catégorie stable par petites $\varinjlim$ , et $f$ un foncteur commutant à ces limites inductives. En d'autres termes, on a}
\begin{enumerate}
\item[\textit{1$^\circ$)}] \textit{$\alpha$ possède les propriétés précédentes, i.e.} $\Ind_\pi(C)$ \textit{est stable par $\varinjlim$ , et le foncteur $\alpha$ commute aux $\varinjlim$ de cardinal $\leqslant\pi$.}
\item[\textit{2$^\circ$)}] \textit{Pour toute $\mathfrak{U}$-catégorie $\mathcal{E}$ stable par petites limites inductives, le foncteur}
\[
\begin{aligned}
g\;&\longmapsto\;g\circ\alpha\\
\uHom_{!}(\Ind_\pi C,\mathcal{E})\;&\longrightarrow\;\uHom_{!\,\pi}(C,\mathcal{E})
\end{aligned}
\tag*{(4.4.16.2)}
\]
\textit{est une équivalence de catégories, où l'indice $!$ (resp. $!\,\pi$) indique la sous-catégorie strictement pleine de la catégorie} $\uHom$
\msp{page 102}
\textit{envisagée, formée des foncteurs qui commutent aux petites $\varinjlim$ (resp. à celles de cardinal $\leqslant\pi$).}
\end{enumerate}

On peut donc dire aussi que $\Ind_\pi(C)$ joue le rôle d'une “enveloppe” stable par petites $\varinjlim$ de la catégorie $C$, quand les $\varinjlim$ de cardinal $\leqslant\pi$, lesquelles existent déjà dans $C$, sont respectées par les plongement de $C$ dans son enveloppe. C'est un résultat nullement tautologique qu'il suffit pour cela d'ajouter à $C$ des $\varinjlim$ \textit{grandes devant} $\pi$ d'objets de $C$, et que plus précisément, la catégorie enveloppe est réalisée par $\Ind_\pi C$ (ou, au choix et moins pléthoriquement, par $\Ind'_\pi C\subset C^\wedge$). % [sic] les plongement -> le plongement, number agreement

\textsc{Démonstration}. Le 1°) est contenu dans 4.4.13 a) et c). Prouvons donc 2°). Il est évident que $\alpha$ est un foncteur pleinement fidèle dont l'image essentielle engendre strictement
\printedpage{27}
% head: A. Grothendieck : Dérivateurs, Ch. XVIII, Édition des Universités de Montpellier II et Paris VII
% foot: Version 17 décembre 2025
$\mathcal{M}=\Ind_\pi C$, d'où résulte que le foncteur (4.4.16.2) est pleinement fidèle, étant induit par le foncteur pleinement fidèle
\[
\alpha^*_{\mathcal{E}}\;:\;\uHom(\mathcal{M},\mathcal{E})\;\longrightarrow\;\uHom(C,\mathcal{E})
\]
\msp{page 103}
sur les sous-catégorie pleines envisagées. Il reste donc à montrer que (4.4.16.2) est essentiellement surjectif. Soit donc $f:C\longrightarrow\mathcal{E}$ dans $\uHom_{!\,\pi}(C,\mathcal{E})$, considérons son extension canonique % [sic] les sous-catégorie pleines -> les sous-catégories pleines, number agreement
\[
\left\{
\begin{aligned}
&Lf\;:\;\Ind(C)\;\longrightarrow\;\mathcal{E}\\
&Lf((X_i)_{i\in I})\;=\;\varinjlim_{i\in I}f(X_i)
\end{aligned}
\right.
\tag*{(4.4.16.3)}
\]
(cf. SGA 4 I 8.7.3), et soit
\[
\psi(f)\;=\;Lf|\Ind_\pi(C)\;. \tag*{(4.4.16.4)}
\]
Rappelons que $Lf$ commute aux petites $\varinjlim$ filtrantes, mais comme le foncteur d'inclusion
\[
\Ind_\pi C\;\longrightarrow\;\Ind(C) \tag*{(4.4.16.5)}
\]
ne commute pas à ces $\varinjlim$ , il n'est pas clair a priori que $\psi(f)$ commute auxdites limites, ni a fortiori qu'il commute aux petites $\varinjlim$ sans plus. Mais on a vu que (4.4.16.5) commute aux $\varinjlim$ \textit{grandes devant} $\pi$, ce qui est contenu dans le corollaire 4.4.15 (vu que les limites inductives \textit{filtrantes} dans $\Ind(C)$ sont les limites \textit{dans} $C^\wedge$) : dire que $\Ind'_\pi(C)\longrightarrow\Ind'(C)$ commute à un certain type de $\varinjlim$ filtrantes, revient à dire qu'il en est ainsi du composé avec $\Ind'_\pi(C)\hookrightarrow C^\wedge$, ce qui est établi par 4.4.15. De ceci résulte que $\psi(f)$ \textit{commute aux} $\varinjlim$ \textit{grandes devant} $\pi$. Pour établir que $g=\psi(f)$ commute
\msp{page 104}
aux petites $\varinjlim$ , il reste donc à voir qu'il commute aux $\varinjlim$ de cardinal $\leqslant\pi$, en vertu de 4.4.14.1 (marge page 98). Pour l'établir, on utilise le théorème 4.4.13 b, en représentant un foncteur $\varphi:J\longrightarrow\Ind_\pi(C)$ (où $\card\Fl J\leqslant\pi$) sous forme indicielle par
\[
\Phi\;:\;I\times J\;\longrightarrow\;C\;, \tag*{(4.4.16.6)}
\]
$I$ étant un ensemble ordonné grand devant $\pi$. On a donc
\[
\varphi(j)\;=\;\varinjlim_{i\in I}{}_{\Ind_\pi C}\,\Phi(i,j)\;,
\]
et comme $g=\psi(f)$ commute aux $\varinjlim$ de type $I$ par ce qu'on a vu, on aura
\[
g(\varphi(j))\;=\;\varinjlim_{i\in I}{}_{(\mathcal{E})}\;g(\Phi(i,j))\;,
\]
d'où
\[
\varinjlim_{j\in J}{}_{(\mathcal{E})}\,g(\varphi(j))\;=\;\varinjlim_{j\in J}\;\varinjlim_{i\in I}\;\underbrace{g(\Phi(i,j))}_{=\,f(\Phi(i,j))}\qquad\text{(limites dans $\mathcal{E}$)}\;. \tag*{$(*)$}
\]
\printedpage{28}
% head: A. Grothendieck : Dérivateurs, Ch. XVIII, Édition des Universités de Montpellier II et Paris VII
% foot: Version 17 décembre 2025

D'autre part, on a
\[
\begin{aligned}
\varinjlim_{j\in J}{}_{\Ind_\pi}\,\varphi(j)\;&\simeq\;\varinjlim_{j\in J}\;\varinjlim_{i\in I}\;\Phi(i,j)\\
&\simeq\;\varinjlim_{i\in I}\;\underbrace{\varinjlim_{j\in J}\;\Phi(i,j)}_{\substack{\text{peut être calculé}\\ \text{dans $C$, soit}\\ \Psi(i)\,\in\,\Ob C}}\;,
\end{aligned}
\]
donc
\[
\varinjlim_{j\in J}{}_{\Ind_\pi}\,\varphi(j)\simeq\varinjlim_{i\in I}\;\Psi(i)\;, \tag*{$(**)$}
\]
\msp{page 105}
d'où en appliquant $g$, qui commute aux $\varinjlim$ de type $I$, et se rappelant maintenant que $g|C\simeq f$,
\[
g(\,\varinjlim_{J}{}_{\Ind_\pi C}\,\varphi(j))\;\simeq\;\varinjlim_{i\in I}\;f(\Psi(i))\;. \tag*{$(*{*}*)$}
\]
On a d'autre part
\[
\Psi(i)\;\overset{\text{déf}}{=}\;\varinjlim_{j\in J}\;\Phi(i,j)\;,\qquad\qquad \varinjlim\;\text{ dans } C\;,
\]
et comme $f$ commute aux $\varinjlim$ de type $J$ (car $\card\Fl J\leqslant\pi$), on en conclut
\[
f(\Psi(i))\;\simeq\;\varinjlim_{j\in J}\;f(\Phi(i,j))\;,
\]
donc par $(*{*}*)$
\[
g(\,\varinjlim_{j\in J}\;\varphi(j))\;\simeq\;\varinjlim_{i\in I}\;\varinjlim_{j\in J}\;f(\Phi(i,j))\;. \tag*{$(*{*}{*}*)$}
\]
En composant avec $(*)$ et utilisant à nouveau la commutativité de $\varinjlim\limits_{I}$ à $\varinjlim\limits_{J}$ , on trouve donc
\[
g(\,\varinjlim\;\varphi(j))\;\overset{\sim}{\longleftarrow}\;\varinjlim\;g(\varphi(j))\;,
\]
ce qui prouve que $g$ commute aux limites inductives de type $J$, et achève de prouver que $g$ est \texttt{[dans]} $\uHom_{!}(\Ind_\pi(C),\mathcal{E})$. Comme $g|C\simeq f$, on en conclut que le foncteur (4.4.16.2) est essentiellement surjectif, ce qui achève de prouver 4.4.16. De plus, on a construit un foncteur quasi-inverse explicite, $\psi$,
\msp{page 106}
s'insérant dans le diagramme commutatif
\[
\begin{tikzcd}[column sep=large, row sep=large]
\uHom_{!\,\pi}(C,\mathcal{E}) \arrow[r, "\psi", "(\sim)"'] \arrow[d, hook, "\text{incl. pl. fid.}"] & \uHom_{!}(\Ind_\pi(C),\mathcal{E}) \arrow[r] \arrow[d, hook, dotted, "\text{pl. fid.}"] & \uHom_{!\,\pi}(\Ind_\pi(C),\mathcal{E}) \\
\uHom(C,\mathcal{E}) \arrow[r, "L", "\substack{\sim\\ \text{(cf. SGA 4 I 8.7.3)}}"'] & \uHom_{!\,\mathrm{filt}}(\Ind(C),\mathcal{E}) \arrow[ur, "R\ \text{(restriction)}"'] &
\end{tikzcd}
\qquad.
\tag*{(4.4.16.6)} % [sic] (4.4.16.6) -> (4.4.16.7)?, the label (4.4.16.6) is already used for Phi on the preceding page
\]

\end{document}
